% -------------------------------------------------------------------------
% WBS ID: RES-NUM-IBC
% Deliverable: Discrete Initial, Boundary and Interface Treatments
% Status: Draft complete
% Evidence status: Regional and local boundary research specification complete; coupling evidence partial
% Review status: Not reviewed
% Last updated: 2026-07-18
% Dependencies: RES-MOD-IBC, RES-DAT-GEO, RES-DAT-SRC, RES-MOD-MUL
% Downstream interfaces: RES-VER-TST, Software boundary/interface implementation
% -------------------------------------------------------------------------

\section{RES-NUM-IBC --- Discrete Initial, Boundary and Interface Treatments}
\label{sec:res-num-ibc}

\subsection{Purpose and Scope}
\label{sec:res-num-ibc-purpose}

% TODO: Define what this Deliverable covers and explicitly excludes.

\subsection{Research Questions}
\label{sec:res-num-ibc-research-questions}

% TODO: State the questions that must be answered before the Deliverable
% can be accepted.

\subsection{Terminology and Definitions}
\label{sec:res-num-ibc-definitions}

% TODO: Define the technical terminology, symbols, quantities and
% classifications used in this Deliverable.

\subsection{Literature Evidence}
\label{sec:res-num-ibc-literature-evidence}

% TODO: Record evidence from peer-reviewed papers, authoritative datasets,
% standards, technical reports and primary documentation.
%
% Do not create a detached literature-summary list. Explain how each source
% supports, contradicts or limits a project decision.

\subsubsection{Takase et al. (2016)}
\label{sec:res-num-ibc-evidence-takase-2016}

\textcite{takase2016hybrid} demonstrated the coupling of independently meshed two-dimensional shallow-water and three-dimensional Navier--Stokes domains through multiple-point constraints and transmission relations. The method enforced compatibility of velocity and pressure despite differences in mesh dimensionality and interface node arrangement.

The paper is retained as evidence for the interface requirements rather than as a numerical formulation to reproduce directly. The present project requires an interface treatment compatible with a finite-volume regional solver, independent three-dimensional discretisation, saved restart states, and both one-way and bidirectional execution.

\subsubsection{Romano et al. (2025)}
\label{sec:res-num-ibc-evidence-romano-2025}

\textcite{romano2025coupling} transferred a time-dependent
three-dimensional tsunami signal into a reduced two-dimensional
propagation model and examined the sensitivity of the resulting model
chain to the coupling-boundary position. The study showed that transfer
accuracy depended on whether the high-fidelity model had already
resolved the nonlinear and geometrically controlled transformations
that the receiving model could not represent.

The paper is retained as evidence for physics-based interface placement,
sequential restart coupling, and preservation of transmitted waveform
content. Its free-surface-based, frequency-domain transfer is treated as
one case-specific implementation rather than as the interface
specification for the present finite-volume shallow-water model.

\subsubsection{Higuera et al. (2013) and Jacobsen et al. (2012)}
\label{sec:res-num-ibc-evidence-wave-boundaries}

\textcite{higuera2013realistic} and
\textcite{jacobsen2012waves2foam} are retained as evidence that
Navier--Stokes/VOF numerical wave tanks require wave-generation and
absorption boundary treatments designed around the intended incident
wave signal. Source role: `3D-WAVEBC`. Project conclusion: local
inlet, outlet and lateral boundary faces shall preserve the incident
history while preventing artificial reflection from contaminating
run-up, overtopping and loading outputs. Limitation: exact relaxation
profiles and numerical coefficients remain case-study verification
choices.

\subsection{Governing Relations and Quantitative Definitions}
\label{sec:res-num-ibc-governing-relations}

% TODO: Record relevant equations, dimensionless groups, empirical models,
% extraction rules or quantitative definitions.
%
% Use align environments for displayed equations.
% Every displayed equation must have a unique \label{eq:...}.
% Do not place punctuation after displayed equations.

\subsection{Discrete Formulation}
\label{sec:res-num-ibc-discrete-formulation}

% TODO: Complete this RES-NUM Domain-specific subsection.

For the regional solver, boundary conditions are imposed as face states
during residual evaluation. Internal and boundary fluxes are:

\begin{align}
  \widehat{\mathbf{F}}_f
  &=
  \widehat{\mathbf{F}}
  \left(
    \mathbf{U}_{i,f},\,
    \mathbf{U}_{j,f},\,
    \mathbf{n}_f
  \right)
  \label{eq:res-num-ibc-regional-internal-face-flux}
  \\
  \widehat{\mathbf{F}}_f
  &=
  \widehat{\mathbf{F}}
  \left(
    \mathbf{U}_{i,f},\,
    \mathbf{U}_{\mathrm{B},f},\,
    \mathbf{n}_f
  \right)
  \label{eq:res-num-ibc-regional-boundary-face-flux}
\end{align}

The baseline regional boundary assignments are:

\begin{itemize}
  \item offshore artificial boundary: characteristic or radiation open
    boundary;
  \item lateral artificial ocean cuts: radiation/open boundary plus
    sponge damping where required;
  \item physical coastline: topography plus wetting--drying, not an
    outer wall boundary;
  \item true wall or resolved structure: reflective boundary state;
  \item debug fallback only: transmissive or zero-gradient state.
\end{itemize}

For a reflective wall, the boundary state preserves depth and reverses
normal velocity:

\begin{align}
  h_{\mathrm{B}}
  &=
  h_{\mathrm{I}}
  \label{eq:res-num-ibc-regional-reflective-depth}
  \\
  u_{n,\mathrm{B}}
  &=
  -u_{n,\mathrm{I}}
  \label{eq:res-num-ibc-regional-reflective-normal-velocity}
  \\
  u_{t,\mathrm{B}}
  &=
  u_{t,\mathrm{I}}
  \label{eq:res-num-ibc-regional-reflective-tangential-velocity}
\end{align}

The two-dimensional-to-three-dimensional transfer shall construct a local boundary or initial state from:
\begin{align}
    \mathcal{D}_{2\rightarrow3} &= \left\{ h,\, \eta,\, \overline{u},\, \overline{v},\, h\overline{u},\, h\overline{v} \right\}_{\Gamma}
    \label{eq:res-num-ibc-two-to-three-data}
\end{align}

The transferred quantities shall be mapped to:

\begin{align}
    \mathcal{D}_{3,\Gamma} &= \left\{ \vec{u}(z),\, p(z),\, \alpha(z) \right\}_{\Gamma}
    \label{eq:res-num-ibc-three-dimensional-interface-data}
\end{align}

The initial lifting method may use a depth-uniform horizontal velocity profile where the interface lies within a hydrostatic, weakly sheared region. Alternative profiles shall remain available where evidence shows that a uniform reconstruction is inadequate. Hydrostatic pressure may be initialised using:

\begin{align}
    p(z) &= \rho g \max \left( \eta-z,\, 0 \right)
    \label{eq:res-num-ibc-hydrostatic-pressure-lifting}
\end{align}

The local three-dimensional solver shall subsequently develop the resolved vertical velocity, shear, separation, and vortical structure.

\subsubsection{Local Three-Dimensional Boundary-Face States}
\label{sec:res-num-ibc-local-boundary-face-states}

For each local boundary face $f$, the discrete boundary state contains:

\begin{align}
  \mathcal{B}_{f}^{\mathrm{3D}}
  &=
  \left\{
    \alpha_f,\,
    \rho_f,\,
    \mu_f,\,
    \mathbf{u}_f,\,
    p_f,\,
    k_f,\,
    \omega_f,\,
    \phi_f
  \right\}
  \label{eq:res-num-ibc-local-boundary-state}
\end{align}

At $\Gamma_{\mathrm{in}}$, the boundary state is reconstructed from
the regional interface history using the inlet free-surface, velocity
and pressure reconstruction rules defined in `RES-MOD-IBC`.
The discrete inlet flux is:

\begin{align}
  \phi_{f,\mathrm{in}}
  &=
  \mathbf{u}_{f,\mathrm{in}}
  \cdot
  \mathbf{S}_f
  \label{eq:res-num-ibc-local-inlet-face-flux}
\end{align}

At $\Gamma_{\mathrm{out}}$ and $\Gamma_{\mathrm{side}}$, the boundary
state shall use an open, radiation or relaxation treatment. The
accepted treatment shall report reflected-wave contamination at the
measurement locations rather than assuming that a zero-gradient state is
non-reflective.

At $\Gamma_{\mathrm{top}}$, the air boundary shall act as an open
atmospheric boundary with a pressure reference compatible with the
pressure-correction equation. At $\Gamma_{\mathrm{T}}$ and
$\Gamma_{\mathrm{B}}$, no penetration is imposed by:

\begin{align}
  \mathbf{u}_{f}
  \cdot
  \mathbf{n}_f
  &=
  0
  \label{eq:res-num-ibc-local-wall-no-penetration}
\end{align}

and the tangential velocity is treated using the selected no-slip or
wall-function discretisation. For the high-Reynolds-number baseline,
scalable wall-function quantities shall be applied consistently to
velocity, $k$, $\omega$ and wall shear on terrain and barrier patches.
The exact wall-function constants and roughness parameters remain open
until the case-study material and validation data are selected.

The three-dimensional-to-two-dimensional restriction operator shall
construct the receiving shallow-water state from the resolved
three-dimensional solution.

The depth shall be obtained from the free-surface and bed elevations:

\begin{align}
    h
    &=
    \eta-z_b
    \label{eq:res-num-ibc-restricted-depth}
\end{align}

The depth-integrated discharges shall be obtained from the resolved
horizontal velocity field:

\begin{align}
    q_x
    &=
    \int_{z_b}^{\eta}
    u\,\dd z
    \label{eq:res-num-ibc-restricted-discharge-x}
    \\
    q_y
    &=
    \int_{z_b}^{\eta}
    v\,\dd z
    \label{eq:res-num-ibc-restricted-discharge-y}
\end{align}

The restricted conservative state is therefore:

\begin{align}
    \vec{U}_{\mathrm{2D},\Gamma}
    &=
    \left[
        h,\,
        q_x,\,
        q_y
    \right]^{\mathsf{T}}
    \label{eq:res-num-ibc-restricted-conservative-state}
\end{align}

A free-surface history alone shall not be treated as sufficient for the
nonlinear shallow-water receiver unless an additional closure supplies
a compatible discharge field. The restriction shall retain the spatial
variation of the three-dimensional solution along the interface rather
than replacing the complete section with one scalar waveform.

\subsection{Conservation Properties}
\label{sec:res-num-ibc-conservation-properties}

% TODO: Complete this RES-NUM Domain-specific subsection.

The normal volume flux transferred across an interface segment shall satisfy:
\begin{align}
    \int_{\Gamma_{\mathrm{3D}}} \vec{u}\cdot\vec{n}\,\dd A &= \int_{\Gamma_{\mathrm{2D}}} h\overline{\vec{u}}\cdot\vec{n}\,\dd s
    \label{eq:res-num-ibc-interface-volume-flux}
\end{align}

The corresponding momentum transfer shall be monitored using a discrete interface balance:

\begin{align}
    \vec{\varepsilon}_{\mathrm{mom}} &= \vec{\Phi}_{\mathrm{mom},\mathrm{3D}} - \vec{\Phi}_{\mathrm{mom},\mathrm{2D}}
    \label{eq:res-num-ibc-interface-momentum-imbalance}
\end{align}

The exact momentum-flux definition shall be consistent with the governing equations and numerical fluxes used by both solvers.

\subsection{Consistency and Formal Accuracy}
\label{sec:res-num-ibc-consistency-formal-accuracy}

% TODO: Complete this RES-NUM Domain-specific subsection.

The interface treatment shall be regarded as a numerical transformation
whose error must be measured separately from the errors of either
solver.

For any transferred scalar or flux quantity $\phi$, the interface error
may be represented as:

\begin{align}
    \varepsilon_{\Gamma,\phi}
    &=
    \phi_{\mathrm{received}}
    -
    \mathcal{T}_{\Gamma}
    \left(
        \phi_{\mathrm{donor}}
    \right)
    \label{eq:res-num-ibc-interface-transfer-error}
\end{align}

where $\mathcal{T}_{\Gamma}$ denotes the complete spatial and temporal
mapping operation.

The mapping shall be assessed for:

\begin{itemize}
    \item amplitude preservation;
    \item phase preservation;
    \item normal mass flux;
    \item depth-integrated momentum flux;
    \item spatial variation along the interface;
    \item relevant wave-frequency content;
    \item wet and dry extent;
    \item accumulated interpolation error.
\end{itemize}

\subsection{Stability Requirements}
\label{sec:res-num-ibc-stability-requirements}

% TODO: Complete this RES-NUM Domain-specific subsection.

\subsection{Mesh and Time-Step Requirements}
\label{sec:res-num-ibc-mesh-time-step-requirements}

% TODO: Complete this RES-NUM Domain-specific subsection.

The interface treatment shall support:
\begin{itemize}
  \item non-conforming spatial meshes;
  \item conservative interpolation or projection;
  \item different two-dimensional and three-dimensional time steps;
  \item temporal interpolation of incident states;
  \item configurable coupling intervals;
  \item subcycling of the local three-dimensional solver;
  \item synchronisation at each bidirectional exchange point.
\end{itemize}
For simultaneous coupling, neither solver shall advance beyond a required exchange time until the interface data from the other solver are available. The software implementation may realise this requirement using a task join, future, condition variable, or equivalent synchronisation primitive.

Saved interface data shall include the physical sampling interval,
solver time, spatial coordinate system, interpolation method, and all
metadata required to reproduce the donor signal deterministically.

The sampling interval $\Delta t_{\Gamma}$ shall satisfy the Nyquist
condition for the highest frequency that materially affects the
quantities of interest:

\begin{align}
    f_{\mathrm{s},\Gamma}
    &=
    \frac{1}{\Delta t_{\Gamma}}
    \geq
    2 f_{\mathrm{max}}
    \label{eq:res-num-ibc-interface-nyquist-condition}
\end{align}

where $f_{\mathrm{max}}$ is the highest physically relevant retained
frequency.

Spectral comparison shall be used as a diagnostic even when both
solvers operate entirely in the time domain. This comparison shall
identify attenuation, amplification, phase distortion, aliasing, or
low-pass filtering introduced by checkpoint storage, interpolation, or
interface reconstruction.

\subsection{Numerical Failure Modes}
\label{sec:res-num-ibc-numerical-failure-modes}

% TODO: Complete this RES-NUM Domain-specific subsection.

Potential interface failure modes include:
\begin{itemize}
  \item loss of mass or momentum during spatial projection;
  \item artificial reflection at the model boundary;
  \item suppression of physical reflected waves in one-way mode;
  \item inconsistent free-surface elevation and pressure;
  \item nonphysical vertical profiles introduced by the lifting operator;
  \item temporal phase error caused by interpolation or coupling lag;
  \item instability caused by explicit bidirectional exchange;
  \item overspecification of the three-dimensional boundary;
  \item wetting or drying mismatch across the interface;
  \item thread or process oversubscription during concurrent comparison runs.
\end{itemize}

The Romano et al.\ implementation demonstrates that a case-specific
transfer may be successful with free-surface elevation as the prescribed
quantity. That choice is not adopted directly for the present nonlinear
shallow-water receiver, whose conservative state also requires
depth-integrated discharge.

Similarly, the paper's frequency-domain mild-slope calculation is not
used as the primary regional model. Nevertheless, the observed
sensitivity to retained frequency content establishes that temporal
downsampling or interpolation may remove physically important
short-period components even in a time-domain coupling implementation.

Potential failure modes derived from this comparison include:

\begin{itemize}
    \item transferring into the reduced model before high-fidelity
          processes have decayed;
    \item omitting the discharge associated with a transferred
          free-surface history;
    \item suppressing spatial variation along a wide interface;
    \item filtering relevant high-frequency components;
    \item contaminating the transferred state with an artificial
          boundary;
    \item placing the interface within a strong geometric contraction,
          expansion, or resonance region;
    \item interpreting agreement at one interface position as evidence
          of interface independence.
\end{itemize}

\subsection{Verification Requirements}
\label{sec:res-num-ibc-verification-requirements}

% TODO: Complete this RES-NUM Domain-specific subsection.

The interface implementation shall be tested using:
\begin{itemize}
  \item still-water preservation;
  \item steady uniform-flow transmission;
  \item solitary-wave transmission;
  \item propagation from two dimensions into three dimensions;
  \item propagation from three dimensions back into two dimensions;
  \item reflection from a controlled obstacle;
  \item non-conforming mesh transfer;
  \item time-step and coupling-interval refinement;
  \item interface-location sensitivity;
  \item one-way replay against simultaneous coupling.
\end{itemize}
Each test shall report free-surface error, phase error, reflected-wave error, mass imbalance, momentum imbalance, and the relevant engineering quantities of interest.

\subsection{Candidate Approaches}
\label{sec:res-num-ibc-candidate-approaches}

% TODO: Define the credible candidate approaches considered.

\subsection{Critical Comparison}
\label{sec:res-num-ibc-critical-comparison}

% TODO: Compare candidates using physically and project-relevant criteria.
% Do not select an approach solely on popularity or implementation ease.

\subsection{Assumptions and Applicability}
\label{sec:res-num-ibc-assumptions}

% TODO: State assumptions and the conditions under which they are valid.

\subsection{Limitations and Failure Conditions}
\label{sec:res-num-ibc-limitations}

% TODO: State where the evidence, model or selected approach ceases to be
% reliable or applicable.

\subsection{Project-Specific Conclusions}
\label{sec:res-num-ibc-conclusions}

% TODO: Convert the evidence into conclusions specific to the Studentship
% project. Do not merely restate literature findings.

\subsection{Selected Approach and Rationale}
\label{sec:res-num-ibc-selected-approach}

% TODO: Record the selected approach only after sufficient evidence exists.
% State the alternatives rejected and the reasons for rejection.

The regional boundary treatment is selected at specification level.
Regional-to-local transfer remains an interface workstream: the required
quantities are identified, but final projection operators, synchronised
coupling residuals and acceptance tolerances are not yet selected.

The local boundary treatment is selected at role level: inlet forcing
from regional histories, open/radiation downstream and lateral
boundaries, open atmospheric top and wall-function-compatible
no-penetration/no-slip treatment on terrain and barrier patches. Exact
relaxation profiles, pressure-reference handling, wall functions and
roughness values remain numerical and validation decisions.

\subsection{Unresolved Decisions}
\label{sec:res-num-ibc-unresolved-decisions}

% TODO: Record unresolved questions, required evidence and decision owners.

Open items are final artificial-boundary placement, domain-specific
sponge settings, source/bathymetry projection after data lock, and
regional-to-local coupling operators. Local open items are inlet
relaxation strategy, side/outlet absorption coefficients, pressure
reference treatment, wall-function constants, roughness parameters and
coupling-interface residuals.

\subsection{Requirements and Handoffs}
\label{sec:res-num-ibc-requirements}

% TODO: State requirements passed to other Research Domains, Software,
% verification activities or later execution work.

The Software boundary implementation shall expose boundary-face state
construction for open/radiation, sponge, reflective and transmissive
states. Verification shall test boundary reflection, sponge damping,
wet/dry coastline behaviour and interface transfer separately.

For the local model, Software shall expose boundary-face state
construction for inlet histories, outlet/side radiation or relaxation,
open top and terrain/barrier wall patches. Verification shall test
boundary reflection and inlet reconstruction before local force or
overtopping results are accepted.

\subsection{Deliverable Completion Record}
\label{sec:res-num-ibc-completion-record}

% TODO: Record whether the Deliverable is:
%
% - Not started
% - In progress
% - Draft complete
% - Reviewed
% - Accepted
%
% Acceptance requires:
%
% - the research questions to be answered;
% - claims to be supported by appropriate evidence;
% - assumptions and limitations to be explicit;
% - the project-specific conclusion to be stated;
% - downstream requirements to be traceable;
% - unresolved decisions to be closed or formally deferred.
